\clearpage
\begin{trapbox}
\textbf{误区一：把“斜率相反”理解成“纵坐标相反”}
\begin{center}\begin{tikzpicture}[font=\small,>=stealth,line join=round]
\draw[gray!45,rounded corners=3pt] (0,0) rectangle (14.2,7.2);
\node at (7.1,6.6) {分别找直线的第二交点，纵坐标中的变动部分相消};
\begin{scope}[shift={(1.4,3.8)},scale=.5]
\draw[->,gray] (-.5,0)--(9.7,0) node[above] {$x$};
\draw[->,gray] (0,-4.8)--(0,3.7) node[left] {$y$};
\draw[teal!65!black,thick,samples=110,domain=-4.3:3.4,variable=\t] plot ({\t*\t/2},\t);
\draw[gray,dashed] (-.2,1)--(4.5,1);
\draw[blue!75!black,very thick] (.5,1)--(2,2);
\draw[orange!85!black,very thick] (.5,1)--(8,-4);
\draw[green!45!black,very thick] (2,2)--(8,-4);
\fill (.5,1) circle(2pt) node[above left=10pt] {$A$};
\fill (2,2) circle(2pt) node[above=12pt] {$M$};
\fill (8,-4) circle(2pt) node[above right=12pt] {$N$};
\end{scope}
\node[anchor=west,align=left] at (8.8,3.8)
{$p=1,\quad t=\frac23$\\[14pt]
$u=\dfrac{2p}{t}-p=2$\\[16pt]
$v=-\dfrac{2p}{t}-p=-4$\\[16pt]
$u+v=-2p=-2$};
\node[text width=13.2cm,align=center] at (7.1,.65) {蓝、橙线分别与曲线在 $A$ 之外再交于 $M,N$；绿色弦斜率为 $-1$。};
\end{tikzpicture}
\end{center}


\textbf{错法：}由 $k_{AM}=-k_{AN}$ 直接写 $u=-v$。

\textbf{纠正：}正确关系是
\[
 \frac{2p}{p+u}=-\frac{2p}{p+v}
 \quad\Longrightarrow\quad \boxed{u+v=-2p},
\]
\textbf{不是} $u+v=0$。若误写为 $0$，甚至会把 $MN$ 错判为竖直线。

\medskip
\textbf{误区二：焦点 $F$ 与定点 $A$ 混淆}

$y^2=2px$ 的焦点为 $F=(\frac p2,0)$，而题给 $A=(\frac p2,p)$。二者横坐标相同，\textbf{纵坐标不同}；焦点 $F$ 通常不在抛物线上。

\medskip
\textbf{误区三：省略斜率存在、两点互异的条件}

若 $u=-p$，则 $A$ 和 $M$ 有相同的横坐标、直线 $AM$ 竖直，$k_{AM}$ 不存在；若 $u=p$，则 $M=A$，$AM$ 不能由两个不同点确定。$N$ 同理。题目给出斜率相反，意味着要先排除这些情形。
\end{trapbox}

\clearpage
\begin{trapbox}
\textbf{误区四：直接把任意两点弦写成有限斜率式}

弦公式
\[
 \boxed{k_{P(u)P(v)}=\frac{2p}{u+v}}
\]
要求两点不同，且 $u+v\ne0$。若 $u+v=0$，两点横坐标相等，弦应为\textbf{竖直线}。不过在本题\textbf{成立的前提}下已推出 $u+v=-2p\ne0$，故 $k_{MN}$ 必然存在。

\medskip
\textbf{误区五：以为 $MN$ 永远是同一条直线}
\begin{center}\begin{tikzpicture}[font=\small,>=stealth,line join=round]
\draw[gray!45,rounded corners=3pt] (0,0) rectangle (14.2,9.4);
\node at (7.1,8.8) {方向相同，位置可以不同：两条弦平行而不重合};
\begin{scope}[shift={(2.5,5.0)},scale=.85]
\draw[->,gray] (-.6,0)--(5.6,0) node[above] {$x$};
\draw[->,gray] (0,-4.8)--(0,3.3) node[left] {$y$};
\draw[teal!65!black,thick,samples=110,domain=-4.5:3.1,variable=\t] plot ({\t*\t/4},\t);
\draw[green!45!black,very thick] (0,0)--(4,-4);
\draw[violet,very thick,dashed] (.25,-1)--(2.25,-3);
\fill (0,0) circle(2pt) node[above left=8pt] {$M$};
\fill (4,-4) circle(2pt) node[above right=12pt] {$N$};
\fill (.25,-1) circle(2pt) node[below left=10pt] {$M'$};
\fill (2.25,-3) circle(2pt) node[below left=9pt] {$N'$};
\fill (1,2) circle(2pt) node[above left=10pt] {$A$};
\end{scope}
\node[anchor=west,align=left] at (8.8,5.1)
{$p=2,\quad A(1,2)$\\[15pt]
$MN:\ y=-x$\\[15pt]
$M'N':\ y=-x-\frac34$\\[15pt]
$u+v=u'+v'=-4$};
\node[text width=13.2cm,align=center] at (7.1,.65) {截距由 $-uv/(2p)$ 决定；纵坐标之和固定，不要求乘积固定。};
\end{tikzpicture}
\end{center}


题目只保证 $k_{MN}=-1$，所以 $MN$ \textbf{恒平行于 $y=-x$}，但截距可随 $M,N$ 改变。用两点的纵坐标记参数，弦方程为
\[
 2px-(u+v)y+uv=0.
\]
代入 $u+v=-2p$ 得
\[
 \boxed{y=-x-\frac{uv}{2p}}.
\]
若 $uv$ 变化，截距也变化。

\medskip
\textbf{误区六：推广时忽略 $a=0$}

一般定点纵坐标为 $a$ 时，$k_{MN}=-p/a$ 仅适用于 $a\ne0$；$a=0$ 对应 $MN$ 竖直，不存在有限斜率。高中题中做推广时必须单独讨论这个边界。
\end{trapbox}
